\documentclass[10pt,a4paper]{amsart}
\usepackage[margin=19mm]{geometry}
\usepackage{amsmath,amssymb,mathtools}
\usepackage[scr=rsfs,cal=euler]{mathalfa}
\usepackage{xfrac}
\usepackage[unicode,hidelinks,bookmarks=false]{hyperref}
\newcommand{\ie}{i.e.\ }
\newcommand{\cf}{cf.\ }
\newcommand{\resp}{respectively\ }
\newcommand{\Subsection}{\subsection}
\providecommand{\nobreakdash}{\nobreak\hbox{-}}
\setlength{\emergencystretch}{2em}
\multlinegap0pt
\usepackage{bm}
\usepackage{bbm}
\usepackage{dsfont}
\usepackage{xspace}
\usepackage{cancel}
\usepackage{mathtools}
\usepackage{amsthm}
\makeatletter
\@ifundefined{theorem}{\newtheorem{theorem}{Theorem}}{}
\@ifundefined{proposition}{\newtheorem{proposition}[theorem]{Proposition}}{}
\makeatother
\def\cC{\mathcal{C}}
\def\cH{\mathcal{H}}
\title[Spanning subhypergraphs with degree constraints]{Spanning subhypergraphs with degree constraints}
\author{Noga Alon, Penny Haxell, Aleksa Milojevi\'c, Jacques Verstra\"ete}
\date{Source: arXiv:2609.13002v2 (CC BY 4.0)}
\begin{document}
\maketitle

\noindent\emph{Source and licence.} Spanning subhypergraphs with degree constraints. Version arXiv:2609.13002v2 (CC BY 4.0), distributed under the Creative Commons Attribution 4.0 International licence, as recorded in the arXiv licence metadata for this exact version and shown on the abstract page https://arxiv.org/abs/2609.13002v2. Full rights notice: licenses/arxiv-2609.13002-CC-BY-4.0.txt.\par\smallskip

\noindent\textit{Editorial scope.} Counted unit: the Proposition whose source label is \texttt{prop:4-shallow} in the article (source0.tex lines 534--536, source label \texttt{prop:4-shallow}), delivered together with the article's own complete original proof of it (source0.tex lines 537--553, one \texttt{proof} environment of 17 lines). No mathematics is written, repaired or re-derived: statement and proof are the article's own text line for line. Required context retained uncounted: prop:5-shallow (lines 450--452). Every definition, display and label of those lines is retained unchanged. Omitted from this package and not redistributed: 1--449, 453--533, 554--556. Nothing inside a retained excerpt is omitted or altered: every retained body is delivered exactly as the article's lines are written, so no omission receipt of any kind accompanies this package. Citations: the retained excerpts carry no citation command at all, so no bibliography entry of the article is transcribed and no reference is redistributed. Reference adaptation: none was needed and none was made; every reference inside the delivered unit resolves inside it, so no unresolved reference or citation is produced. Printed slips kept literal: no printed word, symbol, hypothesis or number of the counted statement or of its proof was altered. Layout and class: the article's own class is \texttt{article} with packages including inputenc, fontenc, graphicx, amsmath, amsthm, amsfonts, xcolor, amssymb, tikz, url, hyperref, breakurl, fullpage, esvect; the wrapper is the lane's pinned \texttt{amsart} editorial wrapper at 10pt on A4, which restates the theorem-like environments the retained text needs and adds the article's own macro definitions that the retained text needs (\texttt{\textbackslash cC}, \texttt{\textbackslash cH}), with extra wrapper packages (bm, bbm, dsfont, xspace, cancel, mathtools). The delivered numbering is the unit's own. Assets: no retained span contains a figure environment, a graphics-inclusion command, a PDF-page inclusion, a table or a TikZ picture; no image file is copied into this package, the package declares no asset dependency and no image file is redistributed with it. Licence: version arXiv:2609.13002v2 of this article is distributed under the Creative Commons Attribution 4.0 International licence, as recorded in the arXiv licence metadata for this exact version and shown on the abstract page https://arxiv.org/abs/2609.13002v2. A full rights notice is delivered at licenses/arxiv-2609.13002-CC-BY-4.0.txt. Characters: every retained excerpt is pure ASCII, so the delivered engine, XeLaTeX with its default Latin Modern text font, reports no missing character.
\begin{proposition}\label{prop:5-shallow}
If $\cH$ is a nonempty regular $3$-uniform hypergraph, then $\cH$ has a spanning subhypergraph with degrees in the set $\{1, 2, 3, 4, 5\}$.
\end{proposition}

\begin{proposition}\label{prop:4-shallow}
If $\cH$ is a nonempty regular $3$-uniform hypergraph, then $\cH$ has a spanning subhypergraph with degrees in the set $\{1, 2, 3, 4\}$.
\end{proposition}

\begin{proof}
We start from the subhypergraph $\cC$ obtained in Proposition~\ref{prop:5-shallow} and remove some edges to reduce its maximum degree by one. Note that we treat $\cC$ as a multihypergraph.

The main observation is that if a vertex $v$ has $\deg_{\cC}v=5$, then we may delete the edge containing $v^{(1)}$ in the matching $M$ (which we included in order to cover $v$) and reduce the degree of $v$ in this way. However, one must be careful in which order the vertices $v$ are chosen, in order to ensure that every vertex is still covered after several edges get deleted in this process.

We call the edge containing $v^{(1)}$ in the matching $M$ the \textit{special edge} of $v$, and we use the same term to refer to the corresponding edge in $\cC$.

Let $D$ be an auxiliary directed graph on the set of vertices of degree $5$. We put a directed edge in $D$ from $u$ to $v$ if $v$ is contained in the special edge of $u$ and $v\neq u$. The goal is to delete some special edges associated to the vertices of $D$, and we will use $D$ to guide the order in which we do this.

The maximum outdegree in $D$ is $2$, and hence there exists a vertex $v$ of indegree at most $2$. We now delete the special edge $e_v$ of the vertex $v$ from $\cC$. 

We claim that even though $e_v$ was deleted, $v$ will still be covered at the end of this procedure. This is simple to see: since $v$ had at most $2$ inedges, at most two additional edges covering $v$ will be deleted. Since $v$ had appeared in $5$ edges of $\cC$, even after all the deletions, $v$ will still be covered. 

We further claim that no other vertex $w$ of $e_v$ is left uncovered. If $w$ never had degree $5$, then the special edge $e_w$ will not be deleted from $\cC$ and $w$ will remain covered. If $w$ had degree $5$ and the special edge $e_w$ was deleted in the past, we have just argued in the previous paragraph that $w$ will remain covered.

We iterate this process on the remaining vertices that are covered at least $5$ times, until all vertices are covered at most $4$ times, which completes the proof.
\end{proof}

\end{document}
